% ===================================================================== Chapter 4
\chapter{Engineering Parameters and Materials}\label{ch:params}

\section{Geometry Parameters}
The geometry of Table~\ref{tab:geometry} was chosen by decision D-008: a bore of 20~mm with $L/D_i$ = 30 gives a turbulent flow that
develops over the first 18 diameters; $D_o/D_i$ = 2 gives a wall whose conduction resistance is small but measurable. For the thickness
study the outer diameter is changed to 36 and 44~mm with the bore unchanged.

\section{Fluid Properties}
Air is modelled with temperature-dependent specific heat, viscosity and conductivity from Incropera and DeWitt, Table A.4 \cite{incropera}
(Table~\ref{tab:air}). A cross-check against a second table \cite{cengel} showed about 3\,\% spread in $k$ and $Pr$, which partly cancel to
less than 2\,\% in $h$. Density is not read from the table, whose values refer to 100~kPa; it is computed from Equation~\eqref{eq:idealgas}
at 101,325~Pa. The CFD inlet density implied by the solved mass flow agrees with the analytical value to 0.006\,\%.

%% generated by gen_tables.py - RE-ANALYSIS 2026
\begin{table}[htbp]
\centering
\small
\caption{Air properties used by the analytical model (at the inlet and outlet bulk temperatures)}
\label{tab:air}
\begin{tabular}{lrrl}
\toprule
Property & 300 K & 368.9 K & Unit \\
\midrule
Density $\rho = p_{op}/(RT)$ & 1.1766 & 0.9569 & kg/m$^3$ \\
Specific heat $c_p$ & 1007 & 1011 & J/(kg\,K) \\
Dynamic viscosity $\mu$ & $1.846\times10^{-5}$ & $2.176\times10^{-5}$ & Pa\,s \\
Thermal conductivity $k$ & 0.02624 & 0.03123 & W/(m\,K) \\
Prandtl number $Pr$ & 0.707 & 0.696 & -- \\
\bottomrule
\end{tabular}
\par\smallskip\parbox{\linewidth}{\footnotesize\raggedright Source: \texttt{02\_Engineering\_Calculations/MATERIAL\_PROPERTIES.md}; $c_p$, $\mu$, $k$, $Pr$ from Incropera et al.\ Table A.4 \cite{incropera}; density computed from the ideal-gas law at 101,325~Pa with $R = 287.058$~J/(kg\,K) because tabulated densities refer to 100~kPa. In Fluent, $c_p$, $\mu$ and $k$ are piecewise-linear over 8 points between 250 and 600~K.}
\end{table}


\section{Inconel 718 Properties}
The solid is age-hardened Inconel 718. Temperature-dependent elastic modulus, conductivity, specific heat and proof strength are taken from
the VDM data sheet \cite{vdm718}; the mean thermal expansion coefficient and density from Special Metals data \cite{htm718}
(Table~\ref{tab:inconel}). Poisson's ratio appears in neither data sheet and is an assumption (0.294, literature range 0.284--0.294). The
data are generic datasheet values, not lot-specific properties.

%% generated by gen_tables.py - RE-ANALYSIS 2026
\begin{table}[htbp]
\centering
\small
\caption{Inconel 718 (age-hardened) properties as a function of temperature}
\label{tab:inconel}
\begin{tabular}{lrrrrr}
\toprule
Property & 20\,\textdegree C & 100\,\textdegree C & 200\,\textdegree C & 300\,\textdegree C & 400\,\textdegree C \\
\midrule
Young's modulus $E$ [GPa] & 204 & 199 & 193 & 187 & 180 \\
Thermal conductivity $k_s$ [W/(m\,K)] & 11.5 & 12.1 & 13.5 & 15.2 & 17.1 \\
Specific heat $c_{p,s}$ [J/(kg\,K)] & 460 & 458 & 468 & 485 & 501 \\
0.2\,\% proof strength $S_y$ [MPa] & 1030 & 1060 & 1040 & 1020 & 1000 \\
\midrule
\multicolumn{6}{@{}l}{Mean CTE from 21\,\textdegree C: 12.8 / 13.3 / 13.9 / 14.2 / 14.8 $\times10^{-6}$/K for means to 93 / 204 / 316 / 427 / 538\,\textdegree C} \\
\multicolumn{6}{@{}l}{Density $\rho_s$ = 8190 kg/m$^3$ (constant); Poisson's ratio $\nu$ = 0.294 [ASSUMED], literature range 0.284--0.294} \\
\bottomrule
\end{tabular}
\par\smallskip\parbox{\linewidth}{\footnotesize\raggedright Sources: VDM Alloy 718 Data Sheet No.~4127 \cite{vdm718} ($E$, $k_s$, $c_{p,s}$, $S_y$); High Temp Metals / Special Metals technical data \cite{htm718} (mean CTE, density). Poisson's ratio is in neither datasheet and is an assumption (A-021). $S_y$ values are typical, not minimum-guaranteed (A-022). In Mechanical the secant CTE is re-referenced from its 70\,\textdegree F datum to $T_{ref} = 300$~K (MPAMOD).}
\end{table}


\section{Temperature-Dependent Properties}
Across the solved temperature range the conductivity of Inconel varies by about 32\,\% and the elastic modulus by about 8\,\%, so constant
properties would change the answer. Fluent uses piecewise-linear $k_s(T)$ and $c_{p,s}(T)$; Mechanical uses $E(T)$ and the secant
$\alpha(T)$, re-referenced from the 70\,\textdegree F datum of the data to $T_{ref}$ = 300~K. The yield strength $S_y(T)$ is used only for
assessment and is interpolated at each node's own temperature. The analytical baseline used single values at the mean temperature
($E$ = 188.1~GPa, $\alpha$ = $13.72\times10^{-6}$/K, $S_y$ = 1023.7~MPa at 281.6\,\textdegree C).

\section{Boundary Conditions}
The CFD boundary conditions are listed in Table~\ref{tab:bc}. The heat flux enters only through \texttt{heated\_outer\_wall}; the interface
temperature is solved, never imposed.

%% generated by gen_tables.py - RE-ANALYSIS 2026
\begin{table}[htbp]
\centering
\small
\caption{CFD boundary conditions (read back from Fluent and verified)}
\label{tab:bc}
\begin{tabular}{>{\raggedright\arraybackslash}p{0.30\linewidth}>{\raggedright\arraybackslash}p{0.22\linewidth}>{\raggedright\arraybackslash}p{0.38\linewidth}}
\toprule
Zone & Type & Value \\
\midrule
\texttt{fluid\_inlet} & velocity inlet & 23.5 m/s normal; 300 K; $I = 4.411$\,\%, $D_h = 0.020$~m \\
\texttt{fluid\_outlet} & pressure outlet & 0 Pa gauge; backflow 300 K; reverse flow prevented \\
\texttt{fluid\_solid\_interface} (+ shadow) & wall, coupled & no-slip; temperature and heat flux solved, never imposed \\
\texttt{heated\_outer\_wall} & wall, heat flux & 8000 W/m$^2$ (the only heat input) \\
\texttt{solid\_inlet\_end}, \texttt{solid\_outlet\_end} & wall, heat flux & 0 W/m$^2$ (adiabatic) \\
\bottomrule
\end{tabular}
\par\smallskip\parbox{\linewidth}{\footnotesize\raggedright Source: \texttt{06\_Fluent\_CFD/FLUENT\_SETUP\_NOTES.md} \S8. A heat-path audit confirmed that the flux is applied exactly once.}
\end{table}


\section{Structural Constraints}
Table~\ref{tab:supports} defines the structural load cases and support scenarios. LC1 is statically determinate, so the duct can expand
freely and any reaction is numerical noise. LC2 holds both complete end faces axially; the three mid-span hoop nodes only remove the rigid-body
rotation. Under LC2 the end faces stay plane and perpendicular to the axis, but the ends are free to translate sideways: in buckling terms
this is a guided (sway) column with $K = 1$. S2 and S3 add lateral restraint at one or both ends. None of the three is a model of a real
flange, whose stiffness is undefined.

%% generated by gen_tables.py - RE-ANALYSIS 2026
\begin{table}[htbp]
\centering
\small
\caption{Structural load cases and support scenarios (idealised boundary conditions)}
\label{tab:supports}
\begin{tabular}{>{\raggedright\arraybackslash}p{0.20\linewidth}>{\raggedright\arraybackslash}p{0.52\linewidth}>{\raggedright\arraybackslash}p{0.18\linewidth}}
\toprule
Case & Constraint definition & Role \\
\midrule
LC1 free expansion & 3 outer-ring nodes on the inlet face ($\theta$ = 0/120/240\textdegree), $U_\theta = U_z = 0$ in the duct cylindrical system; statically determinate & gradient stress only \\
LC2 = S1 & $U_z = 0$ on both complete end faces; $U_\theta = 0$ at 3 mid-span outer nodes; radial free; ends free to sway, end rotation held (guided column, $K = 1$) & official baseline \\
S2 & $U_z = U_\theta = 0$ on every node of both end faces (clamped--clamped, $K = 0.5$) & sensitivity \\
S3 & inlet face $U_z = U_\theta = 0$ (clamped); outlet on a deformable remote point on the axis, $U = 0$, rotations free (pinned, $K = 0.699$) & sensitivity \\
\bottomrule
\end{tabular}
\par\smallskip\parbox{\linewidth}{\footnotesize\raggedright Source: \texttt{MASTER\_PROJECT\_DATA.csv} rows M019--M022. None of these represents a defined real flange; the real end restraint is undefined (open task T-034).}
\end{table}


\section{Assumptions}
The principal assumptions and their final status after all simulations are listed in Table~\ref{tab:assumptions}. Several Section 1
assumptions were superseded by the solved models (for example constant material properties); the full register of 31 entries with its
evidence is kept in \texttt{MASTER\_ASSUMPTIONS.md}.

%% generated by gen_tables.py - RE-ANALYSIS 2026
\begin{table}[htbp]
\centering
\footnotesize
\caption{Principal modelling assumptions and their final status}
\label{tab:assumptions}
\begin{tabular}{l>{\raggedright\arraybackslash}p{0.42\linewidth}>{\raggedright\arraybackslash}p{0.14\linewidth}>{\raggedright\arraybackslash}p{0.25\linewidth}}
\toprule
ID & Assumption & Status & Evidence \\
\midrule
A-004 & steady state & valid & time-invariant boundary conditions \\
A-005 & incompressible ideal gas & valid & $\Delta p$ 0.43\,\% of $p_{op}$; Mach $\le$ 0.097 \\
A-007 & buoyancy neglected & valid & $Ri = 4.4\times10^{-5}$ \\
A-008 & viscous dissipation neglected & valid & $Br = 1.8\times10^{-3}$ \\
A-009 & one-way thermal $\rightarrow$ structural coupling & valid & flow-area change 0.70\,\% \\
A-010 & uniform outer-wall heat flux & assumed & idealisation \\
A-011 & adiabatic end faces & assumed & idealisation \\
A-012 & linear elastic material & valid for statics & utilisation $\le$ 0.645 in all cases \\
A-014 & hydraulically smooth wall & assumed & not tested \\
A-015 & internal radiation neglected & assumed / open & not tested (T-009) \\
A-016 & stress-free temperature 300 K & valid & TREF in every solver input \\
A-020b & inlet turbulence intensity 4.411\,\% & assumed & influence limited to the first diameters \\
A-021 & Poisson's ratio 0.294 & assumed & in neither datasheet \\
A-022 & $S_y(T)$ typical datasheet values & assumed & not minimum-guaranteed \\
A-024 & idealised supports S1/S2/S3 & assumed / open & real restraint undefined (T-034) \\
A-025 & perfectly straight tube, no imperfection & assumed & linear eigenvalue buckling (T-035) \\
A-026 & linear elastic buckling & assumed & no stress--strain data (T-036) \\
\bottomrule
\end{tabular}
\par\smallskip\parbox{\linewidth}{\footnotesize\raggedright Source: \texttt{11\_Final\_Audit/MASTER\_ASSUMPTIONS.md} \S3 (full register A-001 to A-031).}
\end{table}


\section{Property Validity Ranges}
Table~\ref{tab:validity} compares the ranges of the property data and correlations with the ranges actually reached. Every solved case stayed
inside the air and Inconel property tables at every iteration. The air table ends at 600~K, which set the lower velocity and upper heat-flux
limits of the parametric study: the highest air-side temperature, 583.4~K in case Q03, is 16.6~K below the table end.

%% generated by gen_tables.py - RE-ANALYSIS 2026
\begin{table}[htbp]
\centering
\footnotesize
\caption{Validity ranges of property data and correlations compared with the solved ranges}
\label{tab:validity}
\begin{tabular}{>{\raggedright\arraybackslash}p{0.30\linewidth}>{\raggedright\arraybackslash}p{0.26\linewidth}>{\raggedright\arraybackslash}p{0.34\linewidth}}
\toprule
Data / correlation & Validity range & Range reached in this project \\
\midrule
Air $c_p$, $\mu$, $k$ table & 250--600 K & baseline fluid 300.0--553.4 K; highest air-side temperature 583.4 K (Q03) \\
Inconel $k_s$, $c_{p,s}$ table & 293--673 K & baseline solid 423.97--562.13 K (cell centres); highest 591.13 K (Q03) \\
Inconel $E$, $S_y$ table & 20--400\,\textdegree C & mapped field 150.69--289.41\,\textdegree C; highest 318.0\,\textdegree C (Q03) \\
Petukhov friction factor & $3\times10^3 < Re < 5\times10^6$ & $Re$ = 25,545--32,954 \\
Gnielinski Nusselt number & $3\times10^3 < Re < 5\times10^6$, $0.5 < Pr < 2000$ & inside \\
Dittus--Boelter Nusselt number & $Re > 10^4$, small wall-to-bulk $\Delta T$ & wall-to-bulk $\Delta T$ up to 206 K: outside, used only as an upper reference \\
\bottomrule
\end{tabular}
\par\smallskip\parbox{\linewidth}{\footnotesize\raggedright Sources: \texttt{MATERIAL\_PROPERTIES.md}, \texttt{CFD\_BASELINE\_RESULTS.md} \S8, \texttt{PROJECT\_STATE.md} \S18.5 and \S19.6. No property was extrapolated in any solved case.}
\end{table}

